\documentclass{article}
\usepackage{polyglossia}
\setdefaultlanguage{spanish}
\usepackage{fontspec}
\setmainfont{Latin Modern Roman}
\usepackage{amsmath}
\begin{document}
\section{XeLaTeX con polyglossia}
Documento típico de XeLaTeX: \today. Acentos, ñ, «comillas», €, ½.
\[ \sum_{k=1}^{n} k = \frac{n(n+1)}{2} \]
\end{document}
